\documentclass[../main.tex]{subfiles}

\begin{document}
\chapter*{Abstract}
\label{chap:Abstract}

La qualità dei dati è un fattore critico per le performance dei modelli di deep learning, eppure la letteratura contemporanea si concentra quasi esclusivamente sul design delle architetture, trascurando l'impatto sistematico della corruzione e degli errori nei dati. Questa tesi affronta tale lacuna attraverso uno studio empirico rigoroso sulla robustezza dei modelli DL sotto diverse tipologie di errore. L'obiettivo principale è misurare in modo quantitativo e replicabile come cinque categorie di rumore: \textbf{duplicati}, \textbf{etichette inconsistenti}, \textbf{valori mancanti}, \textbf{rumore numerico} e \textbf{outliers} influenzino le performance di modelli di deep learning su due domini distinti: \textbf{dati tabulari} e \textbf{serie temporali}. Il lavoro si appoggia su PuckTrick \cite{agostini2025pucktricklibrarymakingsynthetic}, una libreria Python open-source sviluppata dall'Università degli Studi di Milano-Bicocca, per l'introduzione controllata di errori nei dataset. Originariamente progettata per l'ambito machine learning, è stata estesa e portata su \textbf{Apache Spark} per abilitare l'elaborazione distribuita su dataset di grandi dimensioni. Il contributo della tesi è articolato su quattro fronti: \textbf{(1)} il porting su Spark dei metodi di PuckTrick per i dati tabulari; \textbf{(2)} il design sperimentale basato su multiple run (20 seed diversi per combinazione) con misurazione di intervalli di confidenza al 95\% su entrambi i domini; \textbf{(3)} un'analisi sistematica dei dataset corrotti, che verifica quantitativamente gli effetti di ciascuna tipologia di rumore sulle distribuzioni delle feature e del target; \textbf{(4)} un'analisi empirica comparativa tra architetture DL (MLP per dati tabulari; LSTM per serie temporali) per identificare quali modelli risultano più robusti a diversi tipi di errore. I risultati empirici sul modello \textbf{MLP} mostrano che il \textbf{rumore sulle etichette} (\textit{labels}) è la tipologia di corruzione \textbf{più destabilizzante}: l'F1-score degrada da $\mathbf{0.9750}$ a $\mathbf{0.6166}$ e l'AUC-ROC crolla da $\mathbf{0.9591}$ a $\mathbf{0.5119}$ al \textbf{50\%} di rumore, confermando la tendenza dei modelli DL a memorizzare il rumore di supervisione documentata in letteratura~\cite{zhang2017understandingdeeplearningrequires}. Il rumore \textit{duplicated} risulta invece \textbf{completamente ininfluente} sulle prestazioni. Per le tipologie \textit{missing}, \textit{noise} e \textit{outliers}, l'impatto è contenuto sull'F1-score ma visibile sull'AUC-ROC, con differenze significative tra le feature: \texttt{Oil\_temperature} risulta la \textbf{più vulnerabile}, \texttt{TP3} è protetta dalla sua correlazione perfetta con \texttt{Reservoirs} e \texttt{DV\_pressure} dalla sua natura quasi-binaria. I risultati sul modello \textbf{LSTM} confermano il rumore \textit{labels} come tipologia più destabilizzante: l'AUC-ROC crolla da $\mathbf{0.9638}$ a $\mathbf{0.3511}$ al \textbf{50\%}, scendendo \textbf{sotto la soglia del classificatore casuale} ($0.5$) con un degrado più pronunciato rispetto all'MLP ($0.5119$). Il rumore \textit{duplicated} risulta anche qui ininfluente. Per \textit{missing}, \textit{noise} e \textit{outliers}, \texttt{Oil\_temperature} si mostra più robusta rispetto all'MLP grazie al contesto temporale delle finestre LSTM, mentre \texttt{TP3} risulta leggermente più vulnerabile per la minore efficacia del fallback su \texttt{Reservoirs}.





\end{document}